\subsection{Recollections on bivariant theories}\label{sec:bivariant}

\begin{num}
In this section, we recall the bivariant formalism that automatically comes out of
the six functors formalism for ring spectra, as explained in \cite{Deg16}.\footnote{Recall
that bivariant theories were introduced by Fulton and MacPherson in \cite{FMP}. We refer the 
interested reader to the introduction of \cite{Deg16} for more historical remarks.}

Let $\EE$ be a spectrum (resp. ring spectrum) over $\Sigma$. For any ($\Sigma$-)scheme $S$
with structural map $p:S \rightarrow \Sigma$, we put $\EE_S=p^*\EE$  (resp. equipped with its ring structure)
--- this produces an \emph{absolute (resp. ring) spectrum} over $\base$
in the sense of \cite[Def. 1.1.1]{Deg16} (for $\mathscr T=\SH$).

Recall from \cite{LZ} and \cite[Th. 2.34]{Khan} that the existence of the exceptional functors
$(f_!,f^!)$ on the motivic category $\SH$ has been extended to all finite type morphisms,
instead of separated morphisms of finite type as in \cite{CD3}.
Therefore, one can slightly extend the definitions of \cite[Sec. 1.2]{Deg16}.\footnote{Unlike in \cite[Def. 1.2.2]{Deg16},
we will not call this the \emph{Borel-Moore homology} of $X/S$ with coefficients in $\E$.
This terminology appears a bit unfortunate as $f^!(\E_S)$ is not a dualizing object of $\SH(X)$ in general.}
\end{num}
\begin{df}\label{df:bivariant}
Let $\E$ be a spectrum over $\Sigma$.
We define the \emph{bivariant theory} with coefficients in $\E$ as follows.
Let $f:X \rightarrow S$ be a morphism of finite type and  $(n,m)$ be a couple of integers,
or more generally of Zariski locally constant functions $X \rightarrow \ZZ$.
The bivariant group of $X/S$ in degree $(n,m)$ and coefficients in $\E$ is:
$$
\E_{n,m}(X/S)=[\un_X(m)[n],f^!\E_S].
$$
The natural cohomology associated with this bivariant theory is for any scheme $X$:
\begin{equation}\label{eq:biv->coh}
\E^{n,m}(X)=\E_{-n,-m}(X/X)=[\un_X,\E_X(m)[n]].
\end{equation}
%Given a vector bundle $V$ over $X$, it is also customary to consider
% the cohomology of the Thom space of $V$:
%$$
%\E^{n,m}(\Th(V))=[\Th(V),\E_X(m)[n]].
%$$
\end{df}

\begin{rem} \emph{Twisting conventions}. There are three possible, and useful, choices of indexes for a bivariant (and therefore also for cohomology/homology theories)
associated with (ring) spectra. Apart from the above choice, one can also introduce:
\begin{enumerate}
%\item the above convention.
\item (See \cite[\textsection 1]{DJK}) For any integer $n \in \ZZ$ and any virtual vector bundle $v$ over $X$,
$$
\E_{n}(X/S,v)=[\Th(v)[n],f^!\E_S].
$$
\item For an integer $n \in \ZZ$ and a graded line bundle $(\cL,r)$ over $X$,
$$
\E_{n}(X/S,\cL,r)=[\Th(\cL)(r-1)[n+2r-2],f^!\E_S].
$$
\end{enumerate}
The indexing (1) is the most general one, and covers all the other cases according to
the following rules:
\begin{align*}
\E_{n,m}(X/S)&=\E_{n-2m}(X/S,\tw m), \\
\E_{n}(X/S,\cL,r)&=\E_n(X/S,[\cL] \oplus \tw{r-1}),
\end{align*}
where $[\cL]$ (resp. $\tw n$) denotes the virtual vector bundles associated with $\cL$
(resp. $\AA^n$).
When dealing with $\SL$-oriented theories,
convention (2) is the natural one (see \cite[\textsection 7]{DFJK}).
When dealing with ($\mathrm{GL}$ or $\Sp$)-oriented ring spectra, convention \eqref{eq:biv->coh} is the most relevant,
because of Thom isomorphisms. 

When following the indexing convention (1), note that the associated cohomology theory is of the form
\[
\E^{n}(X,v)=\E_{-n}(X/X,-v)\simeq [\un_X,\E_X\otimes \Th(v)[n]].
\] 
\end{rem}

Thanks to the motivic six functors formalism, bivariant theories enjoy a rich structure
that we will briefly survey for completeness. We start by the functorial properties.\footnote{We state
 these properties taking care about possible twists even though we will be mainly interested in trivial twists
 when dealing with $G$-oriented spectra.}
%\jean{It might be a good idea to make virtual vector bundles strictly functorial with respect to pull-backs.}
%\fred{Indeed, twisted cohomology is not strictly functorial: it only defines a pseudo-functor  because one only gets that $g^*f^*(v) \simeq (fg)^*(v)$ (+the cocyle condition). There are various ways to correct this (for once, any pseudo functor admits a strictifaction by a theorem of Giraud!), but I am not really sure this is necessary. Most of all, as mentioned in the footnote, the interest of the oriented case is to trivialize (i.e. get rid of) twists! Thus I think the consideration of abstract twisting is a bit out of topic here! :)} 
\begin{prop}\label{prop:biv_functorial}
Consider the assumptions of the previous definition.
\begin{enumerate}
\item \emph{Base change map}. For any cartesian square
$$
\xymatrix@=10pt{
Y\ar[r]^-{f'}\ar[d]\ar@{}|\Delta[rd] & X\ar[d] \\
T\ar_-f[r] & S,
}
$$
any virtual vector bundle $v$ over $X$ and any integer $n\in \ZZ$
there exists a morphism $\Delta^*:\E_{n}(X/S,v) \rightarrow \E_{n}(Y/T,(f')^*v)$, functorial
with respect to horizontal composition of cartesian squares. We sometimes write abusively: $f^*=\Delta^*$.
\item \emph{Proper covariance}. For any proper morphism $p:Y \rightarrow X$ of $S$-schemes, any virtual vector bundle $v$ on $X$ and any integer $n\in\ZZ$
there exists a morphism $p_*:\E_{n}(Y/S,p^*v) \rightarrow \E_{n}(X/S,v)$,
corresponding to a covariant functor on the category of $S$-schemes (with proper morphisms as morphisms).
\item \emph{\'Etale contravariance}. For any étale morphism $f:Y \rightarrow X$ of $S$-schemes, any virtual vector bundle $v$ on $X$ and any integer $n\in\ZZ$
there exists a morphism $f^*:\E_{n}(X/S,v) \rightarrow \E_{n}(Y/S,f^*v)$,
corresponding to a contravariant functor on the category of $S$-schemes and étale morphisms.
\end{enumerate}
These structural maps satisfy the following properties:
\begin{enumerate}[resume]
\item \emph{Homotopy invariance}. For any vector bundle $p:E \rightarrow S$, any $S$-scheme $X$ and any virtual vector bundle $v$ over $X$
the base change map $p^*:\E_{n}(X/S,v) \rightarrow \E_{n}(X \times_S E/E,(p')^*v)$ is an isomorphism.
\item \emph{\'Etale invariance}. For any étale morphism $f:T \rightarrow S$,
there exists an isomorphism $f^\sharp:\E_{n}(X/S,v) \rightarrow \E_{n}(X/T,v)$ which is natural
with respect to base change maps, proper covariance and étale contravariance.
\item \emph{Base change formulas}. Base change maps are compatible with proper covariance and étale covariance.
Moreover, for any cartesian square of $S$-schemes
$$
\xymatrix@=10pt{
W\ar^g[r]\ar_q[d] & Y\ar^p[d] \\
V\ar_f[r] & X
}
$$
such that $p$ is proper and $f$ is étale, one has: $f^*p_*=q_*g^*$.
\item \emph{Localization}. For any closed immersion $i:Z \rightarrow X$ of $S$-schemes, 
with complementary open immersion $j:U \rightarrow X$, and any virtual vector bundle $v$ on $X$ there exists a long exact sequence of the form:
$$
\E_{n}(Z/S,i^*v) \xrightarrow{i_*}
\E_{n}(X/S,v) \xrightarrow{j^*}
\E_{n}(U/S,j^*v) \xrightarrow{\partial}
\E_{n-1}(Z/S,i^*v)
$$
which is natural with respect to base change, proper covariance and  \'etale contravariance.
\item \emph{Nisnevich descent}. For any Nisnevich distinguished square\footnote{i.e. the square is cartesian, $p$ \'etale, $j$ open immersion,
and $p$ induces an isomorphism $(V-W)_{red} \rightarrow (X-U)_{red}$} of $S$-schemes 
$$
\xymatrix@=10pt{
W\ar^-l[r]\ar_-q[d] & V\ar^-p[d] \\
U\ar_-j[r] & X
}
$$
and any virtual vector bundle $v$ on $X$ there exists a long exact sequence:
$$
\E_{n}(X/S,v)
\xrightarrow{j^*+p^*} \E_{n}(U/S,j^*v) \oplus \E_{n}(V/S,p^*v)
\xrightarrow{q^*-l^*} \E_{n}(W/S,q^*j^*v)
\rightarrow \E_{n-1}(X/S,v),
$$
functorial with respect to base change maps, proper covariance and étale contravariance.
\item \emph{cdh-descent}. For any cdh-distinguished square\footnote{i.e. the square is cartesian, $p$ proper, $i$ closed immersion,
and $p$ is an isomorphism over $(X-Z)$} of $S$-schemes 
$$
\xymatrix@=10pt{
T\ar_-q[d]\ar^-k[r] & Y\ar^-p[d] \\
Z\ar_-i[r] & X
}
$$
and any virtual vector bundle $v$ on $X$ there exists a long exact sequence
\[
\E_{n}(T/S,q^*i^*v)
\xrightarrow{q_*+k_*} \E_{n}(Z/S,i^*v) \oplus \E_{n}(Y/S,p^*v)
\xrightarrow{i_*-p_*} \E_{n}(X/S,v)
\rightarrow \E_{n-1}(T/S,q^*i^*v)
\]
which is again functorial with respect to base change maps, proper covariance and étale contravariance.
\end{enumerate}
\end{prop}

%\begin{prop}\label{prop:biv_functorial}
%Consider the assumptions of the previous definition.
%\begin{enumerate}
%\item \emph{Base change map}. For any cartesian square
%$$
%\xymatrix@=10pt{
%Y\ar[r]\ar[d]\ar@{}|\Delta[rd] & X\ar[d] \\
%T\ar_f[r] & S
%}
%$$
%there exists a morphism $\Delta^*:\E_{n,m}(X/S) \rightarrow \E_{n,m}(Y/T)$, functorial
%with respect to horizontal composition of cartesian squares. We sometimes write abusively: $f^*=\Delta^*$.
%\item \emph{Proper covariance}. For any proper morphism $p:Y \rightarrow X$ of $S$-schemes,
%there exists a morphism $p_*:\E_{n,m}(Y/S) \rightarrow \E_{n,m}(X/S)$,
%corresponding to a covariant functor.
%\item \emph{Étale contravariance}. For any étale morphism $f:Y \rightarrow X$ of $S$-schemes,
%there exists a morphism $f^*:\E_{n,m}(X/S) \rightarrow \E_{n,m}(Y/S)$,
%corresponding to a contravariant functor.
%\end{enumerate}
%These structural maps satisfy the following properties:
%\begin{enumerate}[resume]
%\item \emph{Homotopy invariance}. For any vector bundle $E \rightarrow S$ and any $S$-scheme $X$,
%the base change map $\E_{n,m}(X/S) \rightarrow \E_{n,m}(X \times_S E/E)$ is an isomorphism.
%\item \emph{Étale invariance}. For any étale morphism $f:T \rightarrow S$,
%there exists an isomorphism $f^\sharp:\E_{n,m}(X/S) \rightarrow \E_{n,m}(X/T)$ which is natural
%with respect to base change maps, proper covariance and étale contravariance.
%\item \emph{Base change formulas}. Base change maps are compatible with proper covariance and étale covariance.
%Moreover, for any cartesian square of $S$-schemes
%$$
%\xymatrix@=10pt{
%W\ar^g[r]\ar_q[d] & Y\ar^p[d] \\
%V\ar_f[r] & X
%}
%$$
%such that $p$ is proper and $f$ is étale, one has: $f^*p_*=q_*g^*$.
%\item \emph{Localization}. For any closed immersion $i:Z \rightarrow X$ of $S$-schemes, 
%with complementary open immersion $j:U \rightarrow X$, there exists a long exact sequence of the form:
%$$
%\E_{n,m}(Z/S) \xrightarrow{i_*}
%\E_{n,m}(X/S) \xrightarrow{j^*}
%\E_{n,m}(U/S) \xrightarrow{\partial_i}
%\E_{n-1,m}(Z/S)
%$$
%which is natural with respect to base change, proper covariance and  \'etale contravariance.
%\item \emph{Nisnevich descent}. For any Nisnevich distinguished square\footnote{i.e. the square is cartesian, $p$ \'etale, $j$ open aimmersion,
%and $p$ induces an isomorphism $(V-W)_{red} \rightarrow (X-U)_{red}$} of $S$-schemes
%$$
%\xymatrix@=10pt{
%W\ar^l[r]\ar_q[d] & V\ar^p[d] \\
%U\ar^j[r] & X
%}
%$$
%there exists a long exact sequence:
%$$
%\E_{n,m}(X/S)
%\xrightarrow{j^*+p^*} \E_{n,m}(U/S) \oplus \E_{n,m}(V/S)
%\xrightarrow{q^*-l^*} \E_{n,m}(W/S)
%\rightarrow \E_{n-1,m}(X/S),
%$$
%functorial with respect to base change maps, proper covariance and étale contravariance.
%\item \emph{cdh-descent}. For any cdh-distinguished square\footnote{i.e. the square is cartesian, $p$ proper, $i$ closed immersion,
%and $p$ is an isomorphism over $(X-Z)$} of $S$-schemes
%$$
%\xymatrix@=10pt{
%T\ar_q[d]\ar^k[r] & Y\ar^p[d] \\
%Z\ar^i[r] & X
%}
%$$
%there exists a long exact sequence:
%$$
%\E_{n,m}(T/S)
%\xrightarrow{q_*+k_*} \E_{n,m}(Z/S) \oplus \E_{n,m}(Y/S)
%\xrightarrow{i_*-p_*} \E_{n,m}(X/S)
%\rightarrow \E_{n-1,m}(T/S)
%$$
%again functorial with respect to base change maps, proper covariance and étale contravariance.
%\end{enumerate}
%\end{prop}
These properties follow from the Grothendieck six functors formalism,
as stated in \cite[Introduction]{CD3}. This is a good exercise for the interested reader
(see \cite{Deg16} after Prop. 1.2.4 for hints).

\begin{rem}\label{rem:biv->cohsup}
As can be guessed from the localization exact sequence, the bivariant theory can be thought of as an extension of
cohomology with support. Indeed, when considering a closed immersion $i:Z \rightarrow X$ and a virtual vector bundle $v$ on $X$, if we view $Z$ as
an $X$-scheme (of finite type!), then one has:
$$
\E^{n}_Z(X,v)=\E_{-n}(Z/X,-i^*v),
$$
where one can define the left hand-side, cohomology of $X$ with support in $Z$, as:
$$
\E^{n}_Z(X,v)=[\Sigma^\infty X/(X-Z),\E_X\otimes \Th(v)[n]].
$$
So bivariant theory of closed immersions is really cohomology with support (see also \cite[Rem. 1.2.5]{Deg16}).
This can shed light on the bivariant formalism.
\end{rem}


\begin{num}\textit{Theories with proper support}.\label{num:biv_support}
It is possible to introduce
two variants of the bivariant theory for a (ring) spectrum $\E$ over $\base$,
a morphism $f:X \rightarrow S$ of finite type, integers $(n,m) \in \ZZ^2$
and a virtual vector bundle $v$ over $X$
(\cite[Def. 2.2.1]{DJK}):
\begin{itemize}
\item Bivariant theory with proper support:
\begin{align*}
\E_{n,m}^c(X/S)&=[\un_S(m)[n],f_!(f^!\E_S)] \\
\E_{n}^c(X/S,v)&=[\un_S[n],f_!(f^!\E_S \otimes \Th(-v))]
\end{align*}
\item Cohomology with proper support:
\begin{align*}
\E^{n,m}_c(X/S)&=[\un_S,f_!(\E_X)(m)[n]] \\
\E_{n}^c(X/S,v)&=[\un_S,f_!(\E_X \otimes \Th(v))[n]].
\end{align*}
\end{itemize}
Using again the six functors formalism, one deduces that they satisfy properties similar
to that of (plain) bivariant theory:
bivariant theory (resp. cohomology) with proper support is covariant (resp. contravariant) 
in $X$ with respect to any morphism (resp. proper morphisms), contravariant (resp. covariant)
in $X$ with respect to \'etale morphisms. The base change map (point (1)) exists for any cartesian square
in cohomology with proper support,
but only when $f$ is smooth for bivariant homology with proper support.

The properties (4) to (9) hold true for these two theories, with the following changes:
\begin{itemize}
\item for (5), one needs to assume that $f$ is \'etale and finite;
\item for (7) and cohomology with proper support, one needs to consider $j_*$ and $i^*$
instead of $i_*$ and $j^*$;
\item for (9) and cohomology with proper support,
the exact sequence  must be formulated in terms of pullbacks.
\end{itemize}
\end{num}

We next explain product structures on (plain) bivariant theories coming from ring spectra. We state the definition and properties for bivariant theories under the convention \eqref{eq:biv->coh}. The reader may consult \cite[\S 2.2.7 (4)]{DJK} for the most general indexing convention. 


\begin{prop}\label{prop:biv_product}
Let $\E$ be a ring spectrum over $\Sigma$ and consider the assumptions of the previous definition.

For any composable morphisms of finite type $Y \xrightarrow g X \xrightarrow f S$,
there exists a product of the form
$$
\E_{n,m}(Y/X)\otimes \E_{s,t}(X/S) \rightarrow \E_{n+s,m+t}(Y/S), (y,x) \mapsto y.x
$$
which satisfies the following properties:
\begin{enumerate}
\item \emph{Associativity}.-- Consider four morphisms of finite type $Z/Y/X/S$,
and a triple $(z,y,x) \in  \E_{**}(Z/Y) \times \E_{**}(Y/X) \times \E_{**}(X/S)$, we have:
$$(z.y).x=z.(y.x).$$
\item \emph{Compatibility with pullbacks}.-- Consider morphisms of finite type $Y/X/S$ and any morphism $f:S' \rightarrow S$.
We let $g:X' \rightarrow X$ be the pullback of $f$ along $X/S$.
Then for any pair $(y,x) \in \E_{**}(Y/X) \times \E_{**}(X/S)$, the following formula holds:
$$
f^*(y.x)=g^*(y).f^*(x).
$$
\item \emph{(First) projection formula}.-- Consider morphisms $Z \xrightarrow f Y \rightarrow X \rightarrow S$ of finite type
such that $f$ is proper. Then for any pair $(z,y) \in  \E_{**}(Z/Y) \times \E_{**}(Y/X)$, one has:
$$
f_*(z.y)=f_*(z).y.
$$
\item \emph{(Second) projection formula}.-- Given a cartesian square in the category
of $S$-schemes of finite type:
$$
\xymatrix@=10pt{
Y'\ar^g[r]\ar[d]\ar@{}|\Delta[rd] & Y\ar[d] \\
X'\ar_f[r] & X,
}
$$
such that $f$ is proper, for any pair $(y,x') \in \E_{**}(Y/X) \times \E_{**}(X'/S)$, one has: 
$$
g_*(\Delta^*(y).x')=y.f_*(x').
$$
\end{enumerate}
\end{prop}
For the construction of the product, we start with maps
$$
y:\un_Y(m)[n] \rightarrow g^!(\E_X), x:\un_X(t)[s] \rightarrow f^!(\E_S)
$$
and then define $y.x$ as the following composite:
\begin{align*}
\un_Y(m+t)[n+s] &\xrightarrow{y(t)[s]} g^!(\E_X)(t)[s] \simeq g^!(\E_X \otimes \un_X(t)[s]) \\
& \xrightarrow{g^!(Id \otimes x)} g^!\big(\E_X \otimes f^!(\E_S)\big)=g^!\big(f^*(\E_S) \otimes f^!(\E_S)\big) \\
& \xrightarrow{Ex^{!*}_\otimes} g^!\big(f^!(\E_S \otimes \E_S)\big) \simeq (fg)^!(\E_S \otimes \E_S) \\
& \xrightarrow{(fg)^!(\mu)} (fg)^!(\E_S).
\end{align*}
Here, the isomorphism on the first line holds as Tate twists are invertible.
Second, the map labeled $Ex_\otimes^{!*}$ stand for the map obtained by adjunction
from canonical map:
$$
f_!\big(f^*(\E_S) \otimes f^!(\E_S)\big) \xrightarrow \sim \E_S \otimes f!f^!(\E_S) \xrightarrow{ad(f_!,f^!)} \E_S \otimes \E_S
$$
where the first map is the projection formula isomorphism and the second one is the unit map of the indicated adjunction.
Finally, $\mu$ is the product map of the ring spectrum $\E$.

The stated properties now follows from the six functors formalism,
again left as an exercise to the reader (hints can be found in the proof of \cite[Prop. 1.2.10]{Deg12}).

\begin{rem}
\begin{enumerate}[leftmargin=*]
\item The above product corresponds to the usual cup-product on cohomology, via formula \eqref{eq:biv->coh}.
More generally, it also induces the cup-product on cohomology with support (recall \Cref{rem:biv->cohsup}).
Given a scheme $X$ and closed subschemes $Z$ and $T$ of $X$, one can indeed recover the usual cup product as follows:
\begin{align*}
\E^{n,m}_T(X) \otimes \E^{s,t}_Z(X)=\E_{-n,-m}(T/X) & \otimes \E_{-s,-t}(Z/X)
\rightarrow \E_{-n,-m}(T/X) \otimes \E_{-s,-t}(Z \times_X T/T) \\
& \rightarrow \E_{-n-s,-m-t}(Z \times_X T/X)=\E^{n+s,m+t}_{Z \times_X T}(X) \\
(\tau,\zeta) &\mapsto \tau.\Delta^*(\zeta)=:\tau \cupp \zeta
\end{align*}
where $\Delta$ stands for the obvious cartesian square.

It is the opportunity to recall that cup-products in $\AA^1$-homotopy are not merely graded commutative, because of the double indexing.
Rather, using the above notation, one gets:
$$
\tau \cupp \zeta=(-1)^{(n-m)(s-t)}\epsilon^{mt}.\zeta \cupp \tau.
$$
This follows from the fact the permutation on the sphere $S^1$ is $(-1)$
and that on the sphere $\GG=\un(1)[1]$ is $\epsilon$ (see e.g. \cite[Lem. 6.1.1]{MorTri}).
\item Note also that given $X/S$,
one gets as a particular case both a left and right action of cohomology on bivariant theory:
\begin{align*}
\E^{s,t}(X) \otimes \E_{n,m}(X/S) & \rightarrow \E_{n-s,m-t}(X/S), \\
\E_{n,m}(X/S) \otimes \E^{s,t}(S) & \rightarrow \E_{n-s,m-t}(X/S).
\end{align*}
\end{enumerate}
\end{rem}

\begin{num}\textit{Fundamental classes}.\label{num:fdl_class}
Let $f:X \rightarrow S$ be a smoothable lci morphism, with virtual tangent bundle $\tau_f$.
According to \cite[Th. 3.3.2, Def. 4.1.3]{DJK}, one associates to $f$ its \emph{fundamental
class}:
$$
\eta_f^\E \in \E_{0,0}(X/S,\tau_f).
$$
By definition, this is the map: $\Th(\tau_f) \rightarrow f^!(\E_S)$ in $\SH(X)$.
Working with $\E$-modules, one obtains by adjunction an $\E$-linear morphism:
\begin{equation}\label{eq:fdl_biv}
\eta_f^\E:\E_X \otimes \Th(\tau_f) \rightarrow f^!(\E_S).
\end{equation}
Note also that this map is simply obtained by evaluating the purity
natural transformation $\pur_f$ of \cite[4.3.1]{DJK} at $\E$.
In particular, if $f$ is smooth, $\eta_f^\E$ is an isomorphism (\cite[Rem. 4.3.8]{DJK}).
Recall more generally the following definition (\cite[Def. 4.3.7]{DJK}):
\end{num}
\begin{df}\label{df:f-pure}
Let $f:X \rightarrow S$ be a smoothable lci map,
 and $\E$ be a ring spectrum over $X$.
 One says that $\E$ is \emph{$f$-pure} if the natural transformation $\eta_f^\E$ defined above is an isomorphism.
\end{df}

\begin{rem}
A spectrum $\E$ over $\ZZ$ is said to be \emph{absolutely pure} (\cite[4.3.11]{DJK}, \cite[1.3.2]{Deg12})
 if for any smoothable morphism $f:X \rightarrow S$ between regular schemes, $\E_S$ is $f$-pure.
 It is known that the K-theory spectrum, the rational motivic Eilenberg MacLane spectrum (\cite[14.4.1]{CD3}),
 the rational sphere spectrum (\cite{DFJK}), and the rational cobordism spectrum are absolute pure.
 It is conjectured that the sphere spectrum, and therefore, the algebraic cobordism spectrum, is absolutely pure.
\end{rem}

\subsection{Orientations and associated virtual Thom classes}

%\fred{1) The arguments here are generalizations of \cite{Deg12, Deg16} as the preceding paragraph.
%I believe one should not write any proof and insert appropriate references for the GL-oriented
%corresponding statement. Except \Cref{prop:virt_thom_G-orient} which require a somewhat specific
%argument for G-torsors?}
%
%\fred{2) Je me suis rendu compte qu'une partie des considérations développées ci-dessous
%étaient déjà expliquées dans \cite{DFJK}, dans le cas particulier des $\SL$-orientations...}

In this section, we will make use of the following terminology.

\begin{df}\label{df:graded_lin_alg_group}
Given a positive integer $d$, a \emph{linear algebraic group of degree $d$} is a $\ZZ$-graded linear algebraic group $G_*$
with a homogeneous embeding $G_* \rightarrow \GL_*$ of degree $d$.
To simplify notation, we simply write $G$ for $G_*$.
\end{df}

\begin{num}\label{num:lin_gps&torsors}
Let us recall the theory of orientations in motivic homotopy theory after Panin and Walter,
following the point of view of \cite[\textsection 2]{DF3}.

Let $G$ be a linear algebraic group of degree $d$.\footnote{The most relevant cases for us are $G_*=\GL_*, \SL_*, \Sp_*$,
corresponding to $d=1, 1, 2$. Recall from \emph{loc. cit.} that one can also consider
the case of $\SL_*^c$. Besides, the orthogonal and spinor cases $ \Orth_*$ and $\Spin_*$ are also of interest.}
We will consider the category $\Tors_G(X)$ of Nisnevich $G$-torsors over a
scheme $X$.\footnote{Recall
this category is a groupoid: every morphism must be an isomorphism.}
In the classical cases, this can be viewed as the category of vector bundles over $X$
with some additional data. In any case, there is a canonical functor
$$
o:\Tors_G(X) \rightarrow \Tors_{\GL}(X)=\VB(X).
$$
The rank $r$ of a $G$-torsor $\cV$ is the rank of the associated vector bundle,
seen as a Zariski locally constant function $r:X \rightarrow \NN$.
We assume there exists a $G$-torsor $\un^G_X$ (called the \emph{unit} $G$-torsor)
such that $o(\un^G_X)=\AA^d_X$.

For convenience, let us recall the definition of $G$-orientation (under the form we use, this definition is due to Panin and Walter,
see \cite[Def. 2.1.3]{DF3})
\end{num}
\begin{df}\label{df:G-orientation}
Let $\E$ be a ring spectrum over $\Sigma$.
A $G$-orientation $\thom$ (or absolute orientation) of $\E$
is the data for all schemes $X$ and all $G$-torsors $\cV$ of rank $r$ of an element
$\thom(\cV) \in \E^{2r,r}(\Th(o(\cV)))$ such that the following properties hold:
\begin{enumerate}
\item \textit{Isomorphisms compatibility}. $\phi^*\thom(\cW)=\thom(\cV)$ for an isomorphism
$\phi:\cV \xrightarrow \sim \cW$.
\item \textit{Pullbacks compatibility}. $f^*\thom(\cV)=\thom(f^{-1}\cV)$ for a morphism of schemes $f:Y \rightarrow X$.
\item \textit{Products compatibility}. $\thom(\cV \oplus \cW)=\thom(\cV) \cdot \thom(\cW)$.
\item \textit{Normalization}. $\thom(\un_X^G)=1_X^\E$ via the identification
$\E^{2d,d}(\Th(\AA^d_X)) \simeq \E^{0,0}(X)$, where $d$ is the degree of $G$.
\end{enumerate}
One also says that $(\E,\thom)$ is a $G$-oriented ring spectrum.
\end{df}

\begin{ex} Recall the following classical examples:
\begin{enumerate}
\item $\GL$-oriented: $\mathbf{H}_{\mathrm{M}}R$ ($R$-linear motivic cohomology), $\KGL$, $\MGL$, any ring spectrum associated
with a mixed Weil theory (\cite{CD2}).
\item $\mathrm{SL}^c$-oriented: $\mathbf{H}_{\mathrm{MW}}R$ ($R$-linear MW-motivic cohomology), 
the ring spectrum $\mathbf{H}M$ associated to any MW-module $M$ over a perfect field (\cite{Feld1, Feld2}),
more generaly any ring spectrum which belong to the perverse ($\delta$-)homotopy heart of $\SH(S)$ (\cite{BD1}).
These examples include the Chow-Witt groups. Other examples include $\KO$, the ring spectrum representing higher Grothendieck-Witt groups, and
$\W=\KO[\eta^{-1}]$ (higher Balmer-Witt groups).
\end{enumerate}
Recall moreover:
\begin{itemize}
\item  If $\E$ is $G$-oriented over $S$, any $\E$-algebra in $\SH(S)$ inherits
an induced $G$-orientation. 
\item We have the following string of implications: $\GL$-oriented $\Rightarrow$ $\SL^c$-oriented $\Rightarrow$ $\SL$-oriented $\Rightarrow$ $\Sp$-oriented.
\end{itemize}
\end{ex}

\begin{num}
Recall that a (commutative) Picard groupoid is a symmetric monoidal category $\mathscr C$
which is a groupoid and such that all objects are invertible for the tensor product.
We will generically denote by $+$ the tensor product of all the groupoids considered
in this paper.

Quillen's Q-construction is a basic tool to construct Picard groupoids (see \cite{Del_det}).
Given a scheme $X$, we let $\uK(X)$ be the Picard groupoid of \emph{virtual vector bundles} over $X$:
this is the groupoid associated with the category $Q\VB(X)$.
Similarly, in the notation of the preceding definition, we let $\uK^G(X)$ be
the groupoid of \emph{virtual $G$-torsors} associated with the category $Q\Tors_G(X)$,
where the exact structure on $\Tors^G(X)$ is induced by the exact structure on $\VB(X)$
(\emph{i.e.} from exactness of sequences of $\cO_X$-modules).
We let $0^G_X$ be the $0$-object of $\uK^G(X)$.
In particular, we get a canonical functor 
\begin{equation}\label{eq:forget_virtual}
o_X:\uK^G(X) \rightarrow \uK(X)
\end{equation}
A morphism of schemes $f:Y \rightarrow X$ induces a pullback functor $f^*:\uK^G(X) \rightarrow \uK^G(Y)$
compatible with the functors $o_X$ and $o_Y$ (see also \Cref{num:fibred_cat&K} for more discussion).
The rank $\rk(v)$ of a virtual $G$-torsor $v$ is the rank of the associated virtual vector bundle.

Using homotopy invariance, one can easily extend the existence of Thom classes
from the previous definition to virtual objects.
\end{num}
\begin{prop}\label{prop:virt_thom_G-orient}\label{prop:thom_virtual}
Let $(\E,\thom)$ be a $G$-oriented ring spectrum over $\Sigma$.
Then for any virtual $G$-torsor $\tilde v$ of rank $r$,
 there exists a unique class $\thom(\tilde v) \in \E^{2r,r}(\Th(o(\tilde v)))$ such that:
\begin{enumerate}
\item For a $G$-torsor $\cV$, $\thom([\cV])=\thom(\cV)$.
\item The Thom class of virtual $G$-torsors satisfies the same properties as for $G$-torsors:
$$
\phi^*\thom(\tilde w)=\thom(\tilde v), f^*\thom(\tilde v)=\thom(f^*\tilde v), \thom(\tilde v+\tilde w)=\thom(\tilde v).\thom(\tilde w), \thom(0^G_X)=1_X^\EE,
$$
where $\phi:\tilde v \rightarrow \tilde w$ is an isomorphism in the first relation.
\end{enumerate}
\end{prop}
\begin{proof}
The proof works as in \cite[proof of 4.1.1]{Riou}. Given an exact sequence of $G$-torsors:
$$
0 \rightarrow \cV' \rightarrow \cV \rightarrow \cV'' \rightarrow 0,
$$
we have to show that one has $\thom(\cV)=\thom(\cV').\thom(\cV'')$ through
the canonical identification $\Th(o(\cV))=\Th(o(\cV')) \otimes \Th(o(\cV''))$ produced in \emph{loc. cit.}
As the $\infty$-category $\SH$ satisfies Nisnevich descent, is it sufficient to check
this over local henselian schemes. In other words, one can assume the $G$-torsors
in the above sequence are trivial, which implies that the sequence is split,
the splitting being $G$-equivariant.
In this latter case, the desired relation follows from property (3) of \Cref{df:G-orientation}.
\end{proof}




\begin{num}\textit{Cohomology of Thom spaces}.
Let $\E$ be a spectrum over $\Sigma$.
One usually extends the definition of the cohomology associated with $\E$ 
(\Cref{df:bivariant}) to Thom spaces. Given a virtual vector bundle $v$ over a scheme $X$,
and $(n,m) \in \uZ(X)^2$, one puts:
$$
\E^{n,m}(\Th(v)):=[\Th(v),\E_X(m)[n]].
$$
More generally, for any virtual vector bundle $w$ on $X$ and any $n\in\uZ(X)$ we set
\[
\E^n(\Th(v),w):=[\Th(v),\E_X\otimes\Th(w)[n]]\simeq [\Th(v-w),\E_X[n]],
\]
so that $\E^{n,m}(\Th(v))=\E^{n-2m}(\Th(v),\langle m\rangle)=\E^{n-2m}(\Th(v-\langle m\rangle)$. Therefore the locally constant integer $m$ is somewhat redundant, but we keep it for the sake of clarity.

The above abelian group is functorial in $v$ with respect to isomorphisms of virtual vector bundles,
thus producing a functor
$$
\uZ(X)^2 \times \uK(X) \rightarrow \Ab, (n,m,v) \mapsto \E^{n,m}(\Th(v)).
$$
If $\E$ has a ring structure, one obtains a product:
$$
\E^{n,m}(\Th(v)) \otimes_\ZZ \E^{s,t}(\Th(w)) \rightarrow \E^{n+s,m+t}(\Th(v+w)).
$$
In other words, $\E^{**}(\Th(*))$ is a $\uZ(X)^2 \times \uK(X)$-graded ring.
Note in particular that, fixing $v$, $\E^{**}(\Th(v))$ is an $\E^{**}(X)$-graded module.
\end{num}

As in classical cases, one deduces from Thom classes the so-called \emph{Thom isomorphisms}.
\begin{prop}
Let $(\E,\thom)$ be a $G$-oriented ring spectrum.
Then for any virtual $G$-torsor $\tilde v$ of rank $n$ over a scheme $X$, the map
$$
\E^{**}(X) \rightarrow \E^{*+2n,*+n}(\Th(\tilde v)), \lambda \mapsto \lambda.\thom(\tilde v)
$$
is an isomorphism. In other words, $\E^{**}(\Th(\tilde v))$ is a free rank one $\E^{**}(X)$-module
with basis given by the singleton $\big(\thom(\tilde v)\big)$.
\end{prop}
\begin{proof}
By the compatibility with product of Thom classes, one reduces to the case where
$\tilde v$ is associated to the class of a $G$-torsor $\cV$. The question is Nisnevich local
in $X$, so we can assume that $\cV$ is a trivial $G$-torsor. By compatibility with isomorphisms
and products, one can assume that $\cV=\un_X^G$, in which case the result follows
from the normalization property. 
\end{proof}

\begin{num}\label{num:E-modules}
Let $(\E,\thom)$ be a $G$-oriented ring spectrum over $\Sigma$.
For a scheme $X$, we can consider the category $\E_X\modd$ of $\E_X$-modules
in the monoidal category $\SH(X)$.\footnote{Beware that this category is only additive, which is sufficient for our needs. See \ref{} to get a theory with more structure.}
Given a virtual $G$-torsor $v$ of rank $r$,
the Thom class is by definition a map (in the homotopy category)
$$
\thom(\tilde v):\Th(o(\tilde v)) \rightarrow \E_X(r)[2r].
$$
By adjunction, one gets a morphism in $\E_X\modd$:
$$
\gamma_{\thom(\tilde v)}:\E_X \otimes \Th(o(\tilde v)) \rightarrow \E_X(r)[2r]
$$
which after forgetting the module structure can be written as the composite
\begin{equation}\label{eq:thomclass}
\E_X \otimes \Th(\tilde v) \xrightarrow{\Id_{\E_X} \otimes \thom(\tilde v)} \E_X \otimes \E_X(r)[2r]
\xrightarrow{\mu_X} \E_X(r)[2r]
\end{equation}
where $\mu_X$ is the multiplication map of the ring spectrum $\E_X$. In other words,
$\gamma_{\thom(\tilde v)}$ is the multiplication by the class $\thom(\tilde v)$ seen internally.
The next result is merely a reformulation of the preceding proposition.
\end{num}
\begin{cor}\label{cor:stable-thom}
Consider the above notation. Then the map $\gamma_{\thom(v)}$ is an isomorphism
of $\E_X$-modules.
\end{cor}

\begin{num}\label{num:virtual_Thom_iso}
One deduces from the previous corollary Thom isomorphisms for the various theories associated
 with a ring spectrum (\Cref{df:bivariant}, \Cref{num:biv_support}). As an example,
 for a finite type $S$-scheme $X$, locally constant integers $(n,m) \in \uZ(X)^2$ and a virtual $G$-torsor $\tilde v$ of rank $r$,
 one gets Thom isomorphisms:
\begin{equation}\label{eq:virtual_Thom_iso}
\begin{split}
\gamma_{\thom(\tilde v)}&:\E_{n,m}(X/S,\tilde v) \simeq \E_{n+2r,m+r}(X/S) \\
\gamma_{\thom(\tilde v)}&:\E^{n,m}(X,\tilde v) \simeq \E^{n-2r,m-r}(X)
\end{split}
\end{equation}
The second case is obviously a particular case of the first one!
 Both isomorphisms are induced by multiplication with another form of the Thom
 class: $\thom(v) \in \E^{2r,r}(X,-v)$, which corresponds to the Thom class of \Cref{prop:thom_virtual}
 via the isomorphism: $\E^{2r,r}(Th(v)) \simeq \E^{2r,r}(X,-v)$
 (which comes from definitions and the invertibility of the Thom space $\Th(v)$).
\end{num}


%We introduce the following practical terminology.
%\begin{df}\label{df:graded_lin_alg_group}
%Given a positive integer $d$, a \emph{linear algebraic group of degree $d$} is a $\ZZ$-graded linear algebra group $G_*$
%with a homogeneous embeding $G_* \rightarrow \GL_*$ of degree $d$.
%To simplify notation, we simply write $G$ for $G_*$.
%\end{df}

%\begin{num}\label{num:lin_gps&torsors}
%Let us recall the theory of orientations in motivic homotopy theory after Panin and Walter,
%following the point of view of \cite[\textsection 2]{DF3}.
%
%Let $G$ be a linear algebraic group of degree $d$.\footnote{The most relevant cases for us are $G_*=\GL_*, \SL_*, \Sp_*$,
%corresponding to $d=1, 1, 2$. Recall from \emph{loc. cit.} that one can also consider
%the case of $\SL_*^c$. Besides, the orthogonal and spinor cases $ \Orth_*$ and $\Spin_*$ are also of interest.}
%We will consider the category $\Tors_G(X)$ of Nisnevich $G$-torsors over a
%scheme $X$.\footnote{Recall
%this category is a groupoid: every morphism must be an isomorphism.}
%In the classical cases, this can be viewed as the category of vector bundles over $X$
%with some additional data. In any case, there is a canonical map:
%$$
%o:\Tors_G(X) \rightarrow \Tors_{\GL}(X)=\VB(X).
%$$
%The rank $r$ of a $G$-torsor $\cV$ is the rank of the associated vector bundle,
%seen as a Zariski locally constant function $r:X \rightarrow \NN$.
%We assume there exists a $G$-torsor $\un^G_X$ (called the \emph{unit} $G$-torsor)
%such that $o(\un^G_X)=\AA^d_X$.
%
%For convenience, let us recall the definition of $G$-orientation (under the form we use, this definition is due to Panin and Walter,
%see \cite[Def. 2.1.3]{DF3})
%\end{num}
%\begin{df}\label{df:G-orientation}
%Let $\E$ be a ring spectrum over $\Sigma$.
%A $G$-orientation $\thom$ (or absolute orientation) of $\E$
%is the data for all scheme $X$ and all $G$-torsor $\cV$ of rank $r$ of an element
%$\thom(\cV) \in \E^{2r,r}(\Th(o(\cV)))$ such that the following properties hold:
%\begin{enumerate}
%\item \textit{Isomorphisms compatibility}. $\phi^*\thom(\cV)=\thom(\cW)$ for an isomorphism
%$\phi:\cV \xrightarrow \sim \cW$.
%\item \textit{Pullbacks compatibility}. $f^*\thom(\cV)=\thom(f^{-1}\cV)$ for a morphism of schemes $f:Y \rightarrow X$.
%\item \textit{Products compatibility}. $\thom(\cV \oplus \cW)=\thom(\cV) \cdot \thom(\cW)$.
%\item \textit{Normalization}. $\thom(\un_X^G)=1_X^\E$ via the identification
%$\E^{2d,d}(\Th(\AA^r_X)) \simeq \E^{0,0}(X)$, where $d$ is the degree of $G$.
%\end{enumerate}
%One also says that $(\E,\thom)$ is a $G$-oriented ring spectrum.
%\end{df}
%
%\begin{ex} Recall the following classical examples:
%\begin{enumerate}
%\item $\GL$-oriented: $\mathbf{H}_{\mathrm{M}}R$ ($R$-linear motivic cohomology), $\KGL$, $\MGL$, any ring spectrum associated
%with a mixed Weil theory (\cite{CD2}).
%\item $\mathrm{SL}^c$-oriented: $\mathbf{H}_{\mathrm{MW}}R$ ($R$-linear MW-motivic cohomology), 
%the ring spectrum $\mathbf{H}M$ associated to any MW-module $M$ over a perfect field (\cite{Feld1, Feld2}),
%more generaly any ring spectrum which belong to the perverse ($\delta$-)homotopy heart of $\SH(S)$ (\cite{BD1}).
%These examples include the Chow-Witt groups. Other examples include $\GW$, the ring spectrum representing higher Grothendieck-Witt groups, and
%$\W=\GW[\eta^{-1}]$ (higher Balmer-Witt groups).
%\end{enumerate}
%Recall moreover:
%\begin{itemize}
%\item  If $\E$ is $G$-oriented over $S$, any $\E$-algebra in $\SH(S)$ inherits
%an induced $G$-orientation. 
%\item We have the following string of implications: $\GL$-oriented $\Rightarrow$ $\SL^c$-oriented $\Rightarrow$ $\SL$-oriented $\Rightarrow$ $\Sp$-oriented.
%\end{itemize}
%\end{ex}

\subsection{Euler classes}

As usual, one can associate to a Thom class an Euler class (e.g. \cite[Def. 2.1.7]{DF3}). We give here a convenient definition for our purposes.

\begin{df}\label{df:Euler_class}
Consider a $G$-oriented ring spectrum $(E,\thom)$ as above.

Given a vector bundle $V$ over a scheme $X$,
 a stable $G$-orientation of $V$ is a pair $(\cV,\omega)$, where $\cV \in \uK^G(X)$
 and $\omega\colon V\simeq o(\cV)$ in $\uK(X)$.
Given an orientation $(\cV,\omega)$ on $V$,
 one defines the ($G$-oriented) \emph{Euler class} of $(V,\cV,\omega)$ as
$$
e(V,\cV,\omega)=p^*(\th(\cV)) \in  \E^{2r,r}(X),
$$ 
where $r=\mathrm{rank}(V)$ and $p$ is the composite
\[
p:X \xrightarrow{s_0} V \rightarrow \Th(V)\xrightarrow{\omega}\Th(o(\cV))
\]
of the zero section and the quotient map. We say that the stable orientations $(\cV,\omega)$ and $(\cW,\psi)$ are isomorphic if there exists an isomorphism $f:\cV\to \cW$ in $\uK^G(X)$ making the diagram
\[
\xymatrix{V\ar[r]^-{\omega}\ar@{=}[d] & o(\cV)\ar[d]^-{o(f)} \\
V\ar[r]_-{\psi} & o(\cW)}
\]
commutative.
\end{df}

\begin{rem}\label{rem:invariantGisom}
It is easy to check that the $G$-oriented Euler class of $(V,\cV,\omega)$ depends only on the isomorphism class
 of the $G$-orientation $(\cV,\omega)$. 
\end{rem}

%\jean{
%\begin{rem}\label{rem:dependsonisom} In general, the Euler class does depend on $\cV$ and on the choice of $\omega$. Indeed, suppose that there are two $G$-torsors $V$ and $V'$ for which there exists an isomorphism of vector bundles $\omega:V'\to V$ which is not an isomorphism of $G$-torsors. Then, $\omega^*(\thom(V))\neq \thom(V')$ in general and the Euler class may be different. For a concrete example, let $A=\RR[x,y,z]/(x^2+y^2+z^2-1)$ and set $X=\Spec(A)$. Let respectively $P:=\ker (A^3\xrightarrow{(x,y,z)}A)$ and $P':=\ker (A^3\xrightarrow{(-x,y,z)}A)$. If $v=(x,y,z)$ and $w=(-x,y,z)$, we get generators of $P$ of the form $f_1=e_1-xv, f_2=e_2-yv$ and $f_3=e_3-zv$ and generators of $P'$ of the form $t_1=e_1+xw, t_2=e_2-yw$ and $t_3=e_3-zw$. We have generators $\omega_1:=zf_1\wedge f_2-yf_1\wedge f_3+zf_2\wedge f_3$ of $\det(P)$ and $\omega_2:=zt_1\wedge t_2-yt_1\wedge t_3-xt_2\wedge t_3$, yielding symplectic forms $\theta_1:\det(P)\to \OO_X$ and $\theta_2:\det(P')\to \OO_X$ defined by $\theta_i(\omega_i)=1$. Now, the isomorphism
%\[
%\begin{pmatrix} -1 & 0 & 0 \\ 0 & 1 & 0 \\ 0 & 0 & 1\end{pmatrix}:A^3\to A^3
%\]
%induces an isomorphism $\omega:P'\to P$ satisfying $\omega(t_1)=-f_1$ and $\omega(t_i)=f_i$ for $i=2,3$. We obtain that $\det\omega:\det(P')\to \det(P)$ has the property that $\det\omega(\omega_2)=-\omega_1$ and therefore that $\omega$ is \emph{not} an isomorphism of $\mathrm{Sp}_2$-torsors. It is known that $\widetilde{\mathrm{CH}}^2(X)\simeq \ZZ$ with $e(P,\theta_1)=2$ (or actually $e(P,\theta_1)=-2$ depending on the choices). Now, $e(P,\theta_1)=e(P',-\theta_2)$ using the isomorphism $\omega$, but $2=e(P',-\theta_2)=-e(P',\theta_2)$ from which it follows that $e(P',\theta_2)=-2$. So, it depends if we choose $(\cV,\omega)=((P,\theta_1),\mathrm{Id})$ or $(\cV,\omega)=((P',\theta_2),\omega)$.
%\end{rem}
%}

%\jean{
%\begin{rem}\label{rem:invariantGisom}
%Note however that if $(\cV,\omega)$ and $(\cV',\omega')$ are two stable $G$-orientations such that $\omega_2\omega_1^{-1}:o(\cV)\to o(\cV')$ is of the form $\omega_2\omega_1^{-1}=o(\psi)$ for some $\psi:\cV\to \cV'$ then $e(V,\cV,\omega)=e(V,\cV',\omega')$. 
%\end{rem}}

\begin{rem}
When $E$ is $\GL$-oriented, the Euler class of a vector bundle $V$ is the top Chern class of $V$
(see eg. \cite[Rem. 2.4.5]{Deg12}). In Chow-Witt groups (of MW-motivic cohomology),
the Euler class of an $\SL^c$-torsor was first considered by Barge and Morel, and the second author
(see \cite{Fasel08a}).
\end{rem}

%\begin{num}
%Recall that a (commutative) Picard groupoid is a symmetric monoidal category $\mathscr C$
%which is a groupoid and such that all objects are invertible for the tensor product.
%We will generically denote by $+$ the tensor product of all the groupoids considered
%in this paper.
%
%Quillen's Q-construction is a basic tool to construct Picard groupoids (see \cite{Del_det}).
%Given a scheme $X$, we let $\uK(X)$ be the Picard groupoid of \emph{virtual vector bundles} over $X$:
%this is the groupoid associated with the category $Q\VB(X)$.
%Similarly, in the notation of the preceding definition, we let $\uK^G(X)$ be
%the groupoid of \emph{virtual $G$-torsors} associated with the category $Q\Tors_G(X)$,
%where the exact structure on $\Tors^G(X)$ is induced by the exact structure on $\VB(X)$
%(\emph{i.e.} from exactness of sequences of $\cO_X$-modules).
%We let $0^G_X$ be the $0$-object of $\uK^G(X)$.
%In particular, we get a canonical functor 
%\begin{equation}\label{eq:forget_virtual}
%o_X:\uK^G(X) \rightarrow \uK(X)
%\end{equation}
%A morphism of schemes $f:Y \rightarrow X$ induces a pullback functor $f^*:\uK^G(X) \rightarrow \uK^G(Y)$
%compatible with the functors $o_X$ and $o_Y$ (see also \Cref{num:fibred_cat&K} for more discussion).
%The rank $\rk(v)$ of a virtual $G$-torsor $v$ is the rank of the associated virtual vector.
%
%Using homotopy invariance, one can easily extend the existence of Thom classes
%from the previous definition to virtual objects.
%\end{num}
%\begin{prop}\label{prop:virt_thom_G-orient}\label{prop:thom_virtual}
%Let $(\E,\thom)$ be a $G$-oriented ring spectrum over $\Sigma$.
%Then for any virtual $G$-torsor $\tilde v$ of rank $r$,
% there exists a unique class $\thom(\tilde v) \in \E^{2r,r}(\Th(o(\tilde v))$ such that:
%\begin{enumerate}
%\item For a $G$-torsor $\cV$, $\thom([\cV])=\thom(\cV)$.
%\item The thom class of virutal $G$-torsors satisfies the same properties than for $G$-torsors:
%$$
%\phi^*\thom(\tilde v)=\thom(\tilde w), f^*\thom(\tilde v)=\thom(f^*\tilde v), \thom(\tilde v+\tilde w)=\thom(\tilde v).\thom(\tilde w), \thom(0^G_X)=1_X^\EE,
%$$
%where $\phi:\tilde v \rightarrow \tilde w$ is an isomorphism in the first relation.
%\end{enumerate}
%\end{prop}
%\begin{proof}
%The proof works as in \cite[proof of 4.1.1]{Riou}. Given an exact sequence of $G$-torsors:
%$$
%0 \rightarrow \cV' \rightarrow \cV \rightarrow \cV'' \rightarrow 0,
%$$
%we have to show that one has $\thom(\cV)=\thom(\cV').\thom(\cV'')$ through
%the canonical identification $\Th(o(\cV))=\Th(o(\cV')) \otimes \Th(o(\cV''))$ produced in \emph{loc. cit.}
%As the $\infty$-category $\SH$ satisfies Nisnevich descent, is it sufficient to check
%this over local henselian schemes. In other words, one can assume the $G$-torsors
%in the above sequence are trivial, which implies that the sequence is split,
%the splitting being $G$-equivariant.
%In this latter case, the desired relation follows from property (3) of \Cref{df:G-orientation}.
%\end{proof}
%
\begin{rem}\label{ex:virtual_Euler_classes}
Note that the definition of Euler classes (\Cref{df:Euler_class})
cannot be extended to virtual $G$-torsors $\tilde v$ over a scheme $X$,
as there is no map $X \rightarrow \Th(\tilde v)$ in general.
\end{rem}

%As in classical cases, one deduces from Thom classes the so-called \emph{Thom isomorphisms}.
%\begin{prop}
%Let $(\E,\thom)$ be a $G$-oriented ring spectrum.
%Then for any virtual $G$-torsor $\tilde v$ of rank $n$ over a scheme $X$, the map
%$$
%\E^{**}(X) \rightarrow \E^{*+2n,*+n}(\Th(\tilde v)), \lambda \mapsto \lambda.\thom(\tilde v)
%$$
%is an isomorphism. In other words, $\E^{**}(\Th(\tilde v))$ is a free rank one $\E^{**}(X)$-module
%with basis given by the singleton $\big(\thom(\tilde v)\big)$.
%\end{prop}
%\begin{proof}
%By the compatibility with product of Thom classes, one reduces to the case where
%$\tilde v$ is associated to the class of a $G$-torsor $\cV$. The question is Nisnevich local
%in $X$, so we can assume that $\cV$ is a trivial $G$-torsor. By compatibility with isomorphisms
%and products, one can assume that $\cV=\un_X^G$, in which case the result follows
%from the normalization property. 
%\end{proof}
%
%\begin{num}\label{num:E-modules}
%Let $(\E,\thom)$ be a $G$-oriented ring spectrum over $\Sigma$.
%For a scheme $X$, we can consider the category $\E_X\modd$ of $\E_X$-modules
%in the monoidal category $\SH(X)$.\footnote{Mind this category is only additive.
%It is sufficient for our needs. See \ref{} to get a theory with more structure.}
%Given a virtual $G$-torsor $v$ of rank $r$,
%the Thom class is by definition a map (in the homotopy category)
%$$
%\thom(\tilde v):\Th(o\tilde v) \rightarrow \E_X(r)[2r].
%$$
%By adjunction, one gets a morphism in $\E_X\modd$:
%$$
%\gamma_{\thom(\tilde v)}:\E_X \otimes \Th(o\tilde v) \rightarrow \E_X(r)[2r]
%$$
%which after forgetting the module structure can be written as the composite
%$$
%\E_X \otimes \Th(\tilde v) \xrightarrow{\Id_{\E_X} \otimes \thom(\tilde v)} \E_X \otimes \E_X(r)[2r]
%\xrightarrow{\mu_X} \E_X(r)[2r]
%$$
%where $\mu_X$ is the multiplication map of the ring spectrum $\E_X$. In other words,
%$\gamma_{\thom(\tilde v)}$ is the multiplication by the class $\thom(\tilde v)$ seen internally.
%The next result is merely a reformulation of the preceding proposition.
%\end{num}
%\begin{cor}\label{cor:stable-thom}
%Consider the notation above. Then the map $\gamma_{\thom(v)}$ is an isomorphism
%of $\E_X$-modules.
%\end{cor}
%
%\begin{num}\label{num:virtual_Thom_iso}
%One deduces from the previous corollary Thom isomorphisms for the various theories associated
% with a ring spectrum (\Cref{df:bivariant}, \Cref{num:biv_support}). As an example,
% for a finite type $S$-scheme $X$, integers $(n,m) \in \ZZ^2$ and a virtual $G$-torsor $\tilde v$,
% one gets Thom isomorphisms:
%\begin{equation}\label{eq:virtual_Thom_iso}
%\begin{split}
%\gamma_{\thom(\tilde v)}&:\E_{n,m}(X/S,v) \simeq \E_{n+2r,m+r}(X/S) \\
%\gamma_{\thom(\tilde v)}&:\E^{n,m}(X,v) \simeq \E_{n+2r,m+r}(X)
%\end{split}
%\end{equation}
%The second case is obviously a particular case of the first one!
% Both isomorphisms are induced by multiplication with another form of the Thom
% class: $\thom(v) \in \E^{2r,r}(X,-v)$, which corresponds to the Thom class of \Cref{prop:thom_virtual}
% via the isomorphism: $\E^{2r,r}(Th(v)) \simeq \E^{2r,r}(X,-v)$
% (which comes from definitions and the invertibility of the Thom space $\Th(v)$).
%\end{num}

%\begin{ex}\label{ex:virtual_Euler_classes}
%As an example, one can recover $G$-oriented Euler classes (\Cref{df:Euler_class}),
% associated to a $G$-oriented ring spectrum $(\E,\thom)$, with more generality.
% Let $V$ be a vector bundle over a scheme $X$.
% Recall that one can associate to $V$ its canonical Euler class (\cite[3.1.2]{DJK})
% $e(V,\E) \in \E^{00}(X,-\tw{V})$ given by the following composite map:
%$$
%\un_X=X \xrightarrow{s_0} V \rightarrow \Th(V) \xrightarrow{\eta_\E \otimes \Id} \E \otimes \Th(V)
%$$
%where $s_0$ is the zero section of $V/X$, the second map is the canonical projection, and $\eta_\E$ is the unit
% map of $\E$.
%
%A \emph{stable $G$-orientation} of $V$ will be a pair consisting of an element $\tilde v \in \uK^G(X)$ 
% and an isomorphism $v \simeq o(\tilde v)$ in $\uK(X)$. Then one can define the $G$-oriented Euler class
% associated with $\tilde v$ as the element $e(\tilde v,\E) \in \E^{2r,r}(X)$ obtained as the image of $e(V,\E)$
%  by the following isomorphism
%$$
%\E^{00}(X,-\tw V) \simeq \E^{00}(X,-o(\tilde v))\xrightarrow{\gamma_{\thom(\tilde v)}} \E^{2r,r}(X).
%$$
%One can check that, when $\tilde v$ is the virtual $G$-torsor associated to a $G$-torsor $\tilde \cV$ whose underlying
% vector bundle is isomorphic to $V$, the preceding Euler class recovers that of \Cref{df:Euler_class}.
%\end{ex}

\subsection{Fibred categories and virtual Thom classes}

\begin{num}\label{num:fibred_cat&K}
There is a more concise way of encoding
a $G$-orientation on a ring spectrum, using the language of \emph{fibred categories}.
We refer the reader to \cite{Vistoli} (see also \cite{Gray}).
 By what is classicaly called the Grothendieck construction,
 a fibred category $\phi:\mathscr S \rightarrow \mathscr B$ (\cite[Def. 3.5]{Vistoli})
 can equivalently be described by a (contravariant) \emph{pseudo-functor} (\cite[Prop. 3.13]{Vistoli}):
$$
\psi:\mathscr B^{\mathrm{op}} \rightarrow \Cat
$$
where $\Cat$ is the category of (small) categories, such that for any object $S$ of $\mathscr B$,
$\psi(S)$ is the \emph{fibre category} of $\phi$ over $S$: \emph{i.e.} the full subcategory of $\mathscr S$ whose objects
$X$ satisfy $\phi(X)=S$. Such a pseudo-functor is uniquely associated to the choice of a \emph{cleavage},
which amounts for any morphism $f:T \rightarrow S$ of a pullback functor
$f^*:\psi(S) \rightarrow \psi(T)$. Recall that the word pseudo-functor refers to the fact that one
only gets isomorphism $(\Id_S)^* \simeq \Id_{\psi(S)}$ and $g^*f^* \simeq (gf)^*$, satisfying
natural compatibility conditions and uniquely determined by the cleavage (see \cite[Def. 3.10]{Vistoli}).

All our fibred categories come with an exclicit choice of cleavage, so that we will consider them
as pseudo-functors. Beware however that a morphism of fibred categories (\cite[Def. 3.6]{Vistoli})
only corresponds to what is called a pseudonatural transformation (see \cite[end of 1.5 and 1.6(c)]{Gray}).
Considering a sub-category $\Cat_0$ of $\Cat$,
a fibred category over $\mathscr B$ in $\Cat_0$ (or simply \emph{$\mathscr B$-fibred $\Cat_0$-category})
will be a fibred category whose pseudo-functor has
the form $\psi_0:\mathscr B \rightarrow \Cat_0$.
In our main example, $\Cat_0$ will the be category of Picard groupoids,
but many other cases can appear (additive, triangulated monoidal).

Consider a graded linear algebraic group $G$ as in \Cref{num:lin_gps&torsors}.
Then $X \mapsto \Tors_G(X)$ is a fibred additive category over $\base$,
whose pullback functor is given by the pullback of $G$-torsors.
One deduces (from the functoriality of the associated Picard groupoid)
that $X \mapsto \uK^G(X)$ is a fibred Picard groupoid over $\base$.
Moreover, the functor \eqref{eq:forget_virtual} induces a morphism
of $\base$-fibred Picard groupoids:
$$
o:\uK^G \rightarrow \uK.
$$

%First, recall that $\uK(X)$ is the fiber
% of a fibered Picard groupoid $\phi_{\uK}:\VB \rightarrow \base$ 
% over the category of schemes (in our case, $\base$).\footnote{For recall
% on fibred category, we refer the reader to \cite{Vistoli}.}
% More precisely, an object of $\VB$ is a pair $(v,X)$ where $X$ is a scheme,
% and $v$ is a virtual vector bundle over $X$.
% The functor $\phi_{\uK}$ associates to $(v,X)$ the scheme $X$.
% There is a natural \emph{cleavage} (cf. \cite[Def. 3.9]{Vistoli}) associated with $\phi_{\uK}$.
% More precisely, there exists a \emph{pseudo-functor} (cf. \cite[Def. 3.10]{Vistoli})
% $\uK:\base \rightarrow \mathscr P$ where $\mathscr P$ is the category of (small)
% Picard groupoids\footnote{We have specialized the framework of \cite{Vistoli}
% from the $2$-category of (small) categories $\mathscr Cat$ to that of
% (small) Picard groupoids $\mathscr P$ as this is our natural framework}:
%\begin{itemize}
%\item for a scheme $X$,
% $\uK(X)$ is the category made of virtual vector bundles over $X$
% (the \emph{fiber} of $\phi_{\uK}$ over $X$)
%\item given a morphism $f:Y \rightarrow X$ of schemes,
% one has a pullback functor $f^*:\uK(X) \rightarrow \uK(Y)$.\footnote{In our concrete case,
% the pseudo-functor is given with the construction of virtual vector bundles.
% Beware that in general, for composable morphisms of schemes $f$ and $g$,
% one only has a natural isomorphism $c_{f,g}:g^* \circ f^* \xrightarrow{\sim} (f \circ g)^*$.
% See \emph{loc. cit.} for more details.}
%\end{itemize}
%The term \emph{fibred Picard groupoid} refers not only to the fact that for all scheme $X$,
% $\uK(X)$ is a Picard groupoid, but also to the fact that for all morphisms of scheme
% $f:Y \rightarrow X$,
% $f^*$ is a symmetric monoidal functor (with respect to the operation $+$ on virtual
% vector bundle).
%
%Similarly, for any $\ZZ.r$-graded algebraic group $G$ as above, $\uK^G$ is a pseudo-functor
% corresponding to a fiberd Picard groupoid on $\base$.


The following definition gives an axiomatization of this situation
(it will be useful in the symplectic and orthogonal cases,
see \Cref{sec:sp_orientaiton}).
\end{num}
\begin{df}\label{df:stable_structure}
Let $d>0$ be an integer.
A \emph{stable structure} of degree $d$ on vector bundles (over $\Sigma$) will be:
\begin{itemize}
\item A pseudo-functor $\uK^\sigma:\base \rightarrow \mathscr P$ 
from  $\base$ to the category of Picard groupoids.
\item A natural transformation of pseudo-functors:
$o:\uK^\sigma \rightarrow \uK$.
\item An integer $d>0$ and a cartesian section\footnote{\emph{i.e.} a collection of objects $\un^\sigma_X
\in \uK^\sigma(X)$ indexed by schemes $X$ in $\base$ and stable under pullbacks} $\un^\sigma$ of $\uK^\sigma$
such that $o(\un_X^\sigma)=\tw d$.
\end{itemize}
As a companion, we slightly extend the first part of \Cref{df:Euler_class}. 
 Let $v$ be a virtual vector bundle $v \in \uK(X)$ (resp. $V$ be a vector bundle over $X$).
 A \emph{$\sigma$-orientation} of $v$ (resp. $V$) is a pair $(\tilde v,\omega)$, where $\tilde v \in \uK^\sigma(X)$
 and $\omega$ is an isomorphism $o(\tilde v) \simeq v$ (resp. $o(\tilde v) \simeq \tw V$).
\end{df}

\begin{ex}\label{ex:stable_G-structure}
Let $G$ be a linear algebraic group of degree $d$.
The fibered Picard groupoid $\uK^G$ and the functor $o$ defines a stable structure $\sigma_G$
on vector bundles, as explained in \Cref{num:lin_gps&torsors},
which we call the \emph{$G$-stable structure}.
The main cases for us are $G=\SL$ ($d=1$) and $G=\Sp$ ($d=2$).
Note that the tautological case $G=\GL$ ($\sigma=\Id$, $d=1$) is also interesting
from the point of view of orientations.
\end{ex}

%\begin{ex}
%Our main examples come from a graded linear algebraic group
%$$
%G=G_* \rightarrow \GL_*,
%$$
%among $G_*=\Sp_{2*}, \SL_*, \Orth_*$.
% Then the category of $G_n$-torsors for the fppf-topology over a scheme $X$ naturally maps to that of 
% vector bundles over $X$. In fact, over a scheme $X$:
%\begin{itemize}
%\item $\Sp_{2n}$-torsors for the fppf-topology, equivalently the Zariski topology,
% corresponds to \emph{symplectic bundles}: that is vector bundles $V$ of rank $2n$
% equipped with a symplectic form
% $\psi:V \otimes V \rightarrow \AA^1_X$ (\emph{i.e.} non-degenerate and anti-symmetric).
% The constant object $\un_X^{\Sp}$ is the $2$-dimensional affine space $\AA^{2}_X$ equipped
% with the symplectic form $\begin{pmatrix} 0 & 1 \\ -1 & 0 \end{pmatrix}$.
%\item $\SL_n$-torsors for the fppf-topology, equivalently the Zariski topology,
% corresponds to what is sometimes called, after Barge and Morel,
% \emph{oriented vector bundles}: vector bundles $V$ of rank $n$ equipped with a trivialization of its determinant
% $\lambda:\det(V) \rightarrow \AA^1_X$. The affine line $\AA^1_X$ admits an obvious orientation,
% and we take this as the constant $\SL_*$-torsor.
%\item $\Orth_n$-torsors for the fppf-topology, equivalently the étale topology (but not the Nisnevich one),
% correspond to the so-called \emph{symmetric bundles}: 
% $V$ of rank $n$ equipped with a non-degenerate symmetric form $\phi:V \otimes V \rightarrow \AA^1_X$.
% The affine line admits an obvious trivial non-degenerate symmetric form and we take
% this as the constant symmetric bundle.
%\end{itemize}
%The first case corresponds to our main example (see the previous section).
% Note that the rank of the constant object is $r=2$ for symplectic bundles,
% while it is $r=1$ in the other cases.
%
%For any scheme $X$, one can define the Picard category of virtual $G_*$-torsors
% over $X$, say $\uK^G(X)$ and one gets a natural map of Picard categories
% (which essentially forgets the extra data):
%$$
%\phi^G:\uK^G(X) \rightarrow \uK(X),
%$$
%obviously compatible with pullbacks.
%\end{ex}
%
%\begin{ex}
%\fred{To be developed: examples come also from Schlichting GW-spectra.}
%\end{ex}

\begin{num}
The stable homotopy category $\SH$ is a monoidal triangulated fibred category over $\base$.
We will consider $\SH^\times$ the underlying Picard groupoid of $\otimes$-invertible objects:
for a scheme $S$, $\SH^\times(S)$ is the monoidal category of spectra which are $\otimes$-invertible,
and whose morphisms are isomorphism is $\SH(S)$.

Given a ring spectrum $\E$ (over $\Sigma$), the fibred structure of $\SH$ induces
a fibred structure on $X \mapsto \E_X\modd$, which is therefore an additive $\base$-fibred category $\E\modd$.
As the ring structure on $\E$ is only considered in the homotopy category,
it is not possible to put a monoidal structure on the latter fibred category. However,
there is a morphism $L_{L_\E}:\SH \rightarrow \E\modd$ of additive fibred categories, induced by the
free $\E$-module functor. Given a scheme $S$, we consider the category $\E_X\modd^\times$
whose objects are the $\otimes$-invertible spectra $\F$, and morphismes from $\F$ to $\F'$ are given 
by the sub-abelian group of
$$
\Hom_{\SH(X)}(\F,\E_X \otimes \F') \simeq \Hom_{\E_X\modd}(\E_X \otimes \F,\E_X \otimes \F')
$$
consisting of isomorphisms of $\E_X$-modules.
There is a natural tensor product $\otimes_\E$ on this category, given on objects by the tensor structure
on $\SH(X)$ and on morphisms by the cup-product:
\begin{align*}
&u:\F \rightarrow \E_X \otimes \G, v:\F' \rightarrow \E_X \otimes \G', \\
& u \otimes_\E v:\F \otimes \F' \xrightarrow{u \otimes v} \E_X \otimes \E_X \otimes \G \otimes \G'
\xrightarrow{\mu_X^\E} \E_X \otimes \G \otimes \G'.
\end{align*}
This defines a $\base$-fibred Picard groupoid $\E\modd^\times$, and the free $\E$-module functor
induces a morphism of fibred Picard groupoids: $\SH^\times \xrightarrow{L_\E} \E\modd^\times$.
\end{num}
\begin{df}\label{df:sigma-orientation}
Consider a stable structure $\sigma$ on vector bundles
and an absolute ring spectrum $\E$.

A \emph{$\sigma$-orientation} on the ring spectrum $\E$ will be the data of a natural isomorphism
$\gamma$ 
% admits \emph{$\sigma$-stable Thom classes}
% if there exists the data of a natural isomorphism $\thom$ 
pictured by the double arrow in the following diagram 
whose objects are pseudo-functors,
and simple arrows are natural transformations of pseudo-functors of fibred Picard groupoids:
$$
\xymatrix@R=6pt@C=30pt{
\uK^\sigma\ar^\sigma[r]\ar_{\rk \circ \sigma}[ddd]\ar@{=>}_{\gamma}[rdd]
& \uK\ar^-{\Th}[r]
& \SH^\times\ar^-{L_\E}[r]
&  \E\modd^\times \\
&&& \\
&&& \\
\uZ\ar_{tw^\E_{\PP^1}}[rrruuu] &&&
}
$$
where $\rk:\uK \rightarrow \uZ$ is the pseudo-functor induced by the rank of a virtual bundle,
and $tw^\E_{\PP^1}$ is the pseudo-functor which to an element $n \in \uZ(X)$, associates
the object
$$
tw_{\PP^1}(n)=\bigoplus_{x \in X^{(0)}} \E(n(x))[2n(x)].
$$
One requires further the following normalization condition:
For any scheme $X$, there exists an isomorphism
$\gamma_{(\un^\sigma_X)} \simeq \Id_{\E(d)[2d]}$ compatible with pullbacks.

In this situation, we call $\gamma$ the \emph{$\sigma$-stable Thom isomorphism} on $\E$.
\end{df}


\begin{num}\label{num:Thom_iso_explicit}
Explicitly, a $\sigma$-orientation is the data for any scheme $X$,
and for every $\tilde v \in \uK^\sigma(X)$, of a \emph{Thom isomorphism} in the
category of $\E$-modules over $X$:
\begin{equation}\label{eq:sigma_Thom_iso_map}
\gamma_{\tilde v}:\E_X \otimes \Th(v) \xrightarrow{\sim} \E_X(r)[2r],
\end{equation}
where $v=o(\tilde v)$ is the associated virtual bundle, and $r$ is the rank of $v$.
In addition, this Thom isomorphism is functorial in $\tilde v$ with respect to isomorphisms,
compatible with base change, and compatible with the addition of $v$.

The isomorphism $\gamma_{\tilde v}$ defines a cohomology class:
\begin{equation}\label{eq:sigma-Thom_class}
\th(\tilde v) \in \E^{2r,r}(\Th(v)) \simeq \Hom_{\E\modd_X}(\E_X \otimes \Th(v),\E_X(r)[2r])
\end{equation}
which will be called the \emph{Thom class} of $\tilde v$ with coefficients in $\E$.
 One readily deduces that multiplication by this Thom class induces isomorphisms:
$$
\E^{*,*}(X) \xrightarrow{\sim} \E^{*+2r,*+r}(\Th(v))
$$
and similarly for the other theories associated with $E$ (as in \Cref{num:virtual_Thom_iso}).
\end{num}

The following result is merely a reformulation of \Cref{cor:stable-thom},
showing that \Cref{df:sigma-orientation} is in fact a generalization of 
\Cref{df:G-orientation}.
\begin{prop}\label{prop:stable_G-orientations}
Let $\E$ be a ring spectrum and $G$ a linear algebraic group of degree $d$ (\Cref{df:graded_lin_alg_group}).
Then the map $\thom \mapsto \gamma_\thom$ of \Cref{cor:stable-thom} gives a bijection between the following structures:
\begin{enumerate}
\item the $G$-orientations $\thom$ on $\E$;
\item the $\sigma_G$-stable Thom isomorphisms $\gamma$ on $\E$ (\Cref{ex:stable_G-structure}).
\end{enumerate}
\end{prop}
\begin{proof}
Indeed, the fact that $\gamma_\thom$ is a natural transformation of pseudo-functors comes from the construction
and the compatibility with pullbacks and products in \Cref{df:G-orientation}. 
\end{proof}

\begin{rem}
Using hermitian K-theory, we will give a finer statement in the symplectic case in \Cref{sec:virt_symplect}.
\end{rem}

%\fred{L'intérêt de la définition précédente vient du résultat suivant: \\
%Une orientation orthogonale (resp. symplectique) correspond à un $\KO$-orientation
%(resp. $\KSp$-orientation). Autrement dit, l'isomorphisme de Thom $\thom(v)$
%associé à un $\Orth$-fibré (resp. $\Sp$-fibré) principal virtuel ne dépend que de la classe
%de $v$ dans $\KO(X)$ (resp. $\KSp(X)$). \\
%Notons en effet, que le pseudo-foncteur $\uK^G \rightarrow \uK$ se factorise à équivalence prêt
%par $\underline{\KO}$ ou $\underline{\KSp}$ respectivement. \\
%Au niveau de la rédaction, je ne sais pas encore comment m'organiser.
%Est-ce que ça vaut le coup de faire la définition abstraite ci-dessus ?
%Elle contient au moins le cas orthogonal, et le cas symplectique, donc je dirais oui...}

\begin{num}\textit{$\sigma$-oriented Euler classes}.\label{num:sigma-Thom&Euler}
As in \Cref{num:virtual_Thom_iso}, one gets Thom isomorphisms \eqref{eq:virtual_Thom_iso}
 associated with an element $\tilde v \in \uK^\sigma(X)$ in the various theories
 (cohomology with/without proper support, bivariant with/without proper support)
 with coefficients in $\E$.

One can therefore extend \Cref{ex:virtual_Euler_classes}.
 Let $V/X$ be a vector bundle of rank $r$,
  with a given $\sigma$-orientation $(\tilde v,\omega)$, where $\tilde v \in \uK^{\sigma}(X)$ and $\omega:V\simeq o(\tilde v)$
  (\Cref{df:sigma-orientation}). Then one defines the associated Euler class as follows:
$$
e(V,\tilde v,\omega,\E)=\gamma_{\thom(\tilde v)}\big(e(V,\E)\big) \in \E^{2r,r}(X).
$$
Note that when $\sigma=\sigma_G$ as above, this definition obviously extend
 \Cref{df:Euler_class}.
 Moreover, as in \Cref{rem:invariantGisom},
 one deduces that once $\omega$ is chosen this $\sigma$-oriented Euler class depends only on the isomorphism
 class of $\tilde v$ in $\uK^\sigma(X)$. In other words,
 one can consider that $\tilde v$ belongs to the abelian group
 $K_0^\sigma(X)=\pi_0(\uK^\sigma(X))$.
\end{num}



%\subsection{$\infty$-categorical orientations}
%\fred{Comme dans le paragraphe précédent mais avec des $\infty$-catégories partout! Cela permet de définir le $G$-cobordisme,
%et même le $\sigma$-cobordisme, et la formulation $\infty$-catégorique des classes de Thom donne automatiquement
%la propriété universelle du $\sigma$-cobordisme.}
%
%\subsection{Cobordism theories}
